% !TeX root = ../papers/nc-p3.tex
\section{片切断だけが欠色を持つ場合の十個の候補台}
\label{sec:klein-single-missing-supports}
\begin{proposition}\label{prop:klein-single-missing-supports}
前節の四元群次数を仮定し、二切断のうち一つだけに欠色がある場合を考える。
正則窓では欠色側の分割は $(0,2,2,2)$、他側の $(3,1,1,1)$ 型は全て排除される。
残る $(2,2,1,1)$ 型でも、次に定めるラベル変換・巡回・反転で場合分けすると、
可能な台は検算コードに記載した十個だけである。
\end{proposition}
\begin{proof}[一つずつの切断と四層の自由グラフ]
文字の平行移動と共通の線型置換で欠色を $b=0$、奇側の降順階数を $e=0,1,2,3$ の順に取る。
標準正則なら他の分割がないことは前節の文字別分解による。
奇切断 $N_e(j,l)$ では $b=0$ 由来の四成分が零で、各々の台は十二ビットで指定される。
偶切断 $M_b(i,k)$, $b\ne0$ は十六ビットで指定できる。
各単独ブロックの台を次の必要十分条件で検査する。
階数1は非空長方形、階数2は前節の零成分の同値閉包である。
階数3では台の二部グラフに三マッチングがあり、完全マッチング数が1でないことが必要十分である。
実際完全マッチングがなければ一般階数3である。二つ以上なら、あるマッチングに属して別のものには属さない成分を取ればよい。その変数について行列式は
一次係数も定数項も恒等零ではない。他の非零値を一般に固定すれば非零の根で行列式だけを消せ、
その三次余因子は非零なので階数3を得る。

先に三つの奇色 $1,2,3$ を列挙し、各偶色について残る四ビットのうち
階数2可能な補完だけを取り、最後に奇色0の台を検査する。こうして
$(3,1,1,1)$ で15936576通り、$(2,2,1,1)$ で115151588通りが残る。
四層各四頂点のグラフに、零行・列から得る自由な座標遷移を加え、
偶位置では対角の四辺も加える。強連結成分内に二つの出辺を持つ頂点があれば、
二つの異なる閉路の長さを反復で合わせて指数増大商を得る。
これを通過する台の数は順に0、2664である。
\end{proof}

\begin{proof}[連続純関係と一文字商]
$(2,2,1,1)$ 型の台を次の作用でまとめる。文字 $0,1$ を固定する $2,3$ の交換、
四位置への共通平行移動、位置2,3だけへの1の同時加算、二位置巡回と
$(a,b,c,d)\mapsto(b,a,d,c)$ による反転を用いる。これらによる2664台の軌道数は74である。
反転は左右の正則性を同時に扱う場合分けに使っており、反対環を添字保持同型とみなす必要はない。

色階数 $r=1,2$ の行集合 $I$ への射影階数の上界を $u(I)$ とする。零行を除き、
二次元の直線の同値閉包において $I$ が占める類の個数と $r$ の小さい方で上から抑えられる。
階数1では非零行が全て同じ行直線を張るので同じ上界が使える。
一行を除く残りで $u<r$ なら、その行だけに支えられた純関係が存在する。
二つの連続位置でこの関係が繋がれば定理\ref{thm:no-consecutive-pure-relations}で排除する。

次に位置 $p$ を座標点 $i$ に縮約する。命題\ref{prop:cut-arrow-growth}の先頭・末尾の商は
自由な後続ラベル集合 $X$、先行ラベル集合 $Z$ がそのまま基底を与える。
残る三文字の総次数 $g$ ごとに中央文字は両端から決まるので、関係を入れる前の次元は $|X||Z|$。
位置 $p+1,p+2$ の射影階数上界を $u_{p+1,b},u_{p+2,b}$ として
\[
 d_i\ \ge\ \sum_{g\in K}\max\left(0,|X||Z|-
     \sum_{c\in Z}u_{p+1,c+g}(X)-
     \sum_{a\in X}u_{p+2,a+g}(Z+a+g)\right).
\]
各和はもう一端の文字を固定して投影する関係行の上界であり、交叉は引く必要がない。
この下界が2以上でも一文字商の反復が指数増大する。
この二検査を通るものは74軌道中10個である。

具体的には台を非負整数として表し、非零語 $(i,j,j+e,i+e)$ が
$12e+4((i+j)-1)+j$ ビットに対応するとする（$i+j\ne0$ と $e$ を整数ラベルとしてビット位置に使う）。
十代表はこの順に十六進で
\[
\begin{array}{llll}
10330ffffff&660666ff9ff&66066ffffff&110120ffffff\\
110330ffffff&111111ffffff&330330ff0ff3&330330ff3ff3\\
330330ffffff&330cc0ffffff&&
\end{array}
\]
となる。\path{computations/verify_klein_single_missing_supports.py} は付属のC++17検算を一時ディレクトリでビルドし、
全ての単独階数検査、四層の有向グラフの推移閉包、74軌道の二つの成長検査を再現する。
\end{proof}
